\section{Diseño de APIs REST: Principios de Fielding y OpenAPI}
\subsection{Los 6 principios de REST y su aplicación en SGROAS}

Roy Fielding definió REST en su tesis doctoral (2000) a partir de restricciones arquitectónicas. A continuación se
analiza cada principio y su grado de cumplimiento en la API de SGROAS.

\begin{enumerate}
    \item \textbf{Client-Server:} separación de responsabilidades. El frontend Angular (cliente) y el backend
    Spring Boot (servidor) son aplicaciones independientes que se comunican únicamente por HTTP. \textbf{Cumplido.}
    \item \textbf{Stateless:} cada petición contiene toda la información necesaria. SGROAS utiliza autenticación
    \textbf{JWT stateless}: el servidor no guarda sesiones (\texttt{SessionCreationPolicy.STATELESS}) y cada
    petición incluye el token en la cookie \texttt{access\_token} (HttpOnly, Secure, SameSite=Strict).
    \textbf{Cumplido.}
    \item \textbf{Cacheable:} las respuestas pueden declarar si son cacheadas. SGROAS usa \texttt{@Cacheable} con
    Redis para los listados paginados (TTL 300 segundos por defecto) y cabeceras \texttt{Cache-Control: no-cache}
    para las respuestas autenticadas. \textbf{Cumplido parcialmente} (la caché es a nivel de servidor).
    \item \textbf{Uniform Interface:} URIs de recursos y verbos HTTP semánticos. La API usa sustantivos para los
    recursos (\texttt{/api/conductores}, \texttt{/api/vehiculos}, \texttt{/api/rutas}, \texttt{/api/asignaciones},
    \texttt{/api/incidentes}, \texttt{/api/usuarios}) y verbos HTTP estándar (GET, POST, PUT, DELETE).
    \textbf{Cumplido.}
    \item \textbf{Layered System:} el cliente no conoce detalles internos. La arquitectura de capas (Controller →
    Service → Repository) y la posibilidad futura de colocar un proxy/nginx delante del backend permiten cumplir
    este principio. \textbf{Cumplido parcialmente} (todavía no hay nginx en producción).
    \item \textbf{Code-On-Demand (opcional):} el servidor puede enviar código ejecutable. No se aplica en SGROAS,
    lo cual es una restricción opcional del estilo. \textbf{No aplicado.}
\end{enumerate}

\subsection{Diseño de URIs REST semánticas}

Las URIs del PFC siguen la regla \textit{recursos como sustantivos y jerarquía por identificación}. Por ejemplo,
\texttt{GET /api/conductores/{id}} identifica un conductor concreto, mientras que una URI como
\texttt{/api/conductores/{id}/asignaciones} expresaría la jerarquía conductor→asignaciones (disponible en el modelo
relacional). Se evitan anti-patrones como verbos en la URI (\texttt{/getConductores}) y la versión de la API se
declara en la cabecera/documentación; el proyecto actualmente no versiona por URL, aunque la especificación
OpenAPI declara la versión 0.9.0-rc. El diseño completo de recursos se muestra en la Tabla~\ref{tab:uris}.

\begin{table}[H]
\centering
\small
\begin{tabular}{@{}cll@{}}
\toprule
\textbf{Recurso} & \textbf{Verbos soportados} & \textbf{Respuesta} \\
\midrule
/api/auth/register, /login, /refresh, /logout & POST & AuthResponse / 204 \\
/api/conductores & GET, POST & Page / 201 \\
/api/conductores/\{id\} & GET, PUT, DELETE & Recurso / 204 \\
/api/usuarios & GET, POST & Page / 201 \\
/api/vehiculos & GET, POST & Page / 201 \\
/api/rutas & GET, POST & Page / 201 \\
/api/asignaciones & GET, POST & Page / 201 \\
/api/incidentes & GET, POST & Page / 201 \\
\bottomrule
\end{tabular}
\caption{Recursos de la API REST de SGROAS (ruta base /api).}
\label{tab:uris}
\end{table}

\subsection{JWT: estructura y flujo de autenticación}

Un JWT (\textit{JSON Web Token}) está compuesto por tres partes separadas por puntos:
\begin{enumerate}
    \item \textbf{Header}: algoritmo de firma (HS256 en SGROAS) y tipo de token.
    \item \textbf{Payload}: \textit{claims} (datos reclamados). En SGROAS se implementaron 7 claims conforme a
    RFC 7519: \texttt{jti}, \texttt{iss}, \texttt{sub}, \texttt{aud}, \texttt{iat}, \texttt{nbf} y \texttt{exp},
    más los claims de negocio \texttt{nombre} y \texttt{rol}.
    \item \textbf{Signature}: firma HMAC-SHA256 calculada con la clave secreta (\texttt{app.jwt.secret}).
\end{enumerate}

El flujo implementado es el siguiente: el usuario se autentica en \texttt{POST /api/auth/login} con sus
credenciales; el servidor valida la contraseña con \texttt{BCryptPasswordEncoder}, genera el JWT con
expiración de 1 hora (\texttt{expiration-ms=3600000}) y un \textit{refresh token} de 7 días
(\texttt{refresh-expiration-ms=604800000}) almacenado en Redis (claves \texttt{refresh:...}); el \textit{access
token} se entrega en una cookie \texttt{HttpOnly + Secure + SameSite=Strict}. El frontend envía la cookie en cada
petición con \texttt{withCredentials}; el filtro \texttt{JwtAuthenticationFilter} lee la cookie o la cabecera
\texttt{Authorization: Bearer}, verifica la firma, comprueba si el \texttt{jti} está en la \textit{blacklist} de
Redis y puebla el \texttt{SecurityContext}.

El problema de la \textbf{revocación} de JWT se resolvió con dos estrategias: \textbf{token blacklist en Redis}
(\texttt{blacklist:<jti>}) para invalidar access tokens al hacer logout y \textbf{refresh token con rotación},
almacenado en Redis con TTL de 7 días. Esta decisión quedó documentada en el \textbf{ADR-003}. En
\texttt{security/JwtService.java} la firma se genera con la librería JJWT 0.12.6:

\begin{lstlisting}[style=javaStyle, label=lst:jwt, caption={JwtService.java — generación del token (fragmento)}]
return Jwts.builder()
        .id(jti)
        .issuer(jwtIssuer)
        .subject(usuario.getEmail())
        .audience().add(jwtAudience).and()
        .issuedAt(ahora)
        .notBefore(ahora)
        .expiration(expiracion)
        .claim("nombre", usuario.getNombre())
        .claim("rol", usuario.getRol().name())
        .signWith(getSigningKey(), Jwts.SIG.HS256)
        .compact();
\end{lstlisting}

\subsection{OpenAPI / Swagger 3.0}

SGROAS documenta su API con \textbf{Springdoc OpenAPI 2.8.6}. La especificación OpenAPI actúa como contrato entre
el equipo de frontend y backend. La configuración (\texttt{OpenApiConfig.java}) define el título
\textbf{``SGROAS API''}, la versión 0.9.0-rc y el esquema de seguridad por cookie. La interfaz interactiva de
Swagger UI queda accesible en \texttt{/api/docs/swagger-ui.html}. Se eligió \textbf{Swagger UI} (interfaz
interactiva que permite ejecutar endpoints) frente a Redoc (documentación estática) por su utilidad en la defensa en
vivo. Adicionalmente, el proyecto incluye la colección Postman exportada en
\texttt{docs/postman/coleccion.json} con 22 peticiones organizadas en carpetas, cumpliendo el criterio de
verificación de la guía.

\begin{lstlisting}[style=javaStyle, label=lst:oa, caption={OpenApiConfig.java — bean OpenAPI con cookie auth (fragmento)}]
@Bean
public OpenAPI openAPI() {
    return new OpenAPI()
            .info(new Info()
                    .title("SGROAS API")
                    .description("Sistema de Gestion de Recursos Operativos...")
                    .version("0.9.0-rc"))
            .addSecurityItem(new SecurityRequirement().addList("cookieAuth"))
            .components(new Components()
                    .addSecuritySchemes("cookieAuth", new SecurityScheme()
                            .type(SecurityScheme.Type.APIKEY)
                            .in(SecurityScheme.In.COOKIE)
                            .name("access_token")));
}
\end{lstlisting}